\section{Application: committing data in a Bitcoin transaction annex}
\label{sec:bitcoin-annex}

\subsection{Why this needs a flat hash, not a tree}

BIP-341 \citep{bip341} defines an optional \emph{annex} field in a Taproot
witness stack: any item present when other conditions are met is treated as
annex data, and its commitment inside the transaction's signature hash is
\begin{equation}
\mathrm{sha\_annex} = \mathrm{SHA256}\big(\mathrm{compact\_size}(\mathrm{len}(\mathrm{annex})) \,\Vert\, \mathrm{annex}\big),
\end{equation}
a single, ordinary, untagged SHA-256 of the length-prefixed annex bytes --
consensus never computes anything tree-shaped over annex content. This is
precisely the setting motivating this paper: if the annex bytes are the
FlatMerkle structured file of Section~\ref{sec:construction} (with the
short compact-size preamble accounted for), then $\mathrm{sha\_annex}$ --
already computed by every consensus-validating node, with zero protocol
changes -- \emph{is} the FlatMerkle digest, and anyone holding the annex
can produce $O(\log N)$ inclusion proofs for any element it contains,
verifiable by anyone who has $\mathrm{sha\_annex}$ (Section~\ref{sec:annex-buys}
below explains why this is \emph{not} the same thing as the transaction's
final signature hash, and what a verifier actually needs).

\subsection{What this buys}
\label{sec:annex-buys}

\begin{itemize}
\item A timestamped, immutable commitment to structured application data,
  inheriting Bitcoin's ordering and immutability guarantees, at the cost of
  ordinary transaction fees for the annex bytes.
\item $O(\log N)$ inclusion proofs verifiable by any third party who holds
  $\mathrm{sha\_annex}$ (a public 32-byte value, obtainable from any node
  that has validated the transaction, or supplied directly alongside the
  proof) and a short proof -- no need to download or re-hash the whole
  annex. How a lightweight verifier obtains an \emph{authenticated}
  $\mathrm{sha\_annex}$ for the confirmed transaction is a separate
  problem, with a real cost trade-off
  (Section~\ref{sec:annex-authentication}).
  $\mathrm{sha\_annex}$ itself is \emph{not} the same object as the
  transaction's final Taproot signature hash: it is one 32-byte component
  folded into the larger \texttt{SigMsg} that the signature hash is
  computed over, not recoverable from the signature hash alone (that would
  require inverting a hash function). A verifier who has only the final
  signature digest, with no other data, cannot check a FlatMerkle proof
  against it directly; what must actually be transmitted alongside the
  proof is $\mathrm{sha\_annex}$ itself.
\item By Theorem~\ref{thm:binding}, a given $\mathrm{sha\_annex}$ value
  cannot be made to accept two different, conflicting values for the same
  index: an accepting proof pins down a unique answer. This does
  \emph{not} mean a publisher gains nothing from an annex that deviates
  from the canonical layout -- Remark~\ref{rem:layout-forgery} exhibits a
  real forgery against exactly that stronger claim -- so a third party who
  merely downloads the annex, hashes it, and assumes every accepting proof
  describes bytes it can literally see at the corresponding offset is
  \emph{not} safe. What is safe: any party willing to hold the full annex
  can run \textsc{ScanForViolations} once ($O(N)$) and, if the annex is
  non-canonical, publish an $O(\log N)$ fraud proof
  (Section~\ref{sec:fraud-proofs}) that any lightweight verifier -- with
  only $\mathrm{sha\_annex}$ -- can check without downloading the annex at
  all.
\end{itemize}

\subsection{Alternatives within a linearly hashed annex}
\label{sec:annex-alternatives}

Under the threat model of Section~\ref{sec:intro}, a design must give an
opponent $O(\log N)$ proofs both of inclusion and of malformation. Three
natural alternatives to FlatMerkle meet the size goal but not both proof
goals (Table~\ref{tab:annex-alternatives}).

\emph{Merkle root as a trailer.} Lay the annex out as $\texttt{0x50}$
header, data, then a 32-byte Merkle root $\rho$ of the data, positioned
so that $\rho \Vert \texttt{0x80} \Vert 0^{*} \Vert \mathrm{len}$ fills
the final SHA-256 block. Given the 32-byte midstate $s$ entering that
block, a verifier checks $\mathrm{sha\_annex} = \C(s, \rho \Vert
\mathit{pad})$ with one call -- binding $\rho$ under
Assumption~\ref{ass:freestart}, the same assumption FlatMerkle uses -- and
then an ordinary Merkle path against $\rho$. Overhead is about 96 bytes,
and inclusion proofs are only 32 bytes longer than a plain Merkle proof.
But a malicious committer can publish $\rho = \mathrm{Merkle}(D')$, where
$D'$ equals the published data except at position $0$. Every inclusion
proof then certifies $D'[0]$, not the published $D[0]$, and exposing this
requires proving the bytes physically at position $0$ of the annex: a
verifier must recompute the linear chain from there to the final block,
$\Theta(N)$ compression calls. (We confirmed the one-call opening against
the reference SHA-256 implementation.)

\emph{A fully serialized Merkle tree.} Publishing every internal node in
the annex, as well as the root, has the same weakness: an inconsistent
node can be placed early in the stream, and proving the physical content
of any position in a linear hash costs $\Theta(N)$ without re-injected
states.

\emph{Root trailer plus an interactive bisection game.} The $\Theta(N)$
malformation check can be turned into $O(\log N)$ rounds of an
interactive game: the committer posts intermediate chain states, the
opponent bisects to a single disputed compression call, and a dispute
mechanism such as BitVMX \citep{lerner2024bitvmx} adjudicates it. This
reaches logarithmic cost per round, but needs $O(\log N)$ on-chain rounds,
the committer's participation, and timeouts to handle a committer who
stops responding. FlatMerkle's fraud proofs (Section~\ref{sec:fraud-proofs})
are a single, non-interactive message that anyone holding the annex can
publish.

\begin{table}[htbp]
\centering
\begin{tabular}{@{}lcccc@{}}
\toprule
Design & Annex size & Inclusion & Malformation & Interactive \\
\midrule
Root trailer            & $P + O(1)$ & $O(\log N)$ & $\Theta(N)$ & no \\
Serialized Merkle tree  & $\approx 2P$ & $O(\log N)$ & $\Theta(N)$ & no \\
Root trailer + bisection & $P + O(1)$ & $O(\log N)$ & $O(\log N)$ per round$^{*}$ & yes \\
FlatMerkle              & $\approx 2P$ & $O(\log N)$ & $O(\log N)$ & no \\
\bottomrule
\end{tabular}
\caption{Committing $P$ bytes ($N$ elements) inside a linearly hashed
annex. ``Inclusion'' and ``Malformation'' give the size of a proof an
opponent can produce and the verifier's work to check it. FlatMerkle
pays roughly $2\times$ in size for non-interactive logarithmic proofs of
both kinds. $^{*}$Over $O(\log N)$ interactive rounds.}
\label{tab:annex-alternatives}
\end{table}

\subsection{What FlatMerkle verification does \emph{not} give you: binding \texorpdfstring{$\mathrm{sha\_annex}$}{sha\_annex} to a specific transaction}
\label{sec:annex-authentication}

Everything above is about \emph{FlatMerkle verification}: given a claimed
$\mathrm{sha\_annex}$ value and a short proof, confirming a specific
element sits at a specific index within the annex that hashes to it. That
is a completely separate question from \emph{Bitcoin authentication}: given
a claimed $\mathrm{sha\_annex}$ value and a specific transaction instance,
confirming that value genuinely is that instance's annex hash, and not an
arbitrary 32 bytes an adversary invented and handed to a lightweight
verifier alongside an otherwise-valid-looking FlatMerkle proof.
FlatMerkle's binding guarantee (Theorem~\ref{thm:binding}) says nothing
about this: it is a statement about the relationship between a digest and
the proofs it accepts, not about how a verifier learns which digest
belongs to which transaction. A verifier who accepts \emph{any} 32 bytes as
``the'' $\mathrm{sha\_annex}$ for a given transaction, with no further
check, can be handed a proof for a completely different, adversary-chosen
annex and never know the difference.

A subtlety in identifying ``a given transaction'' at all: the annex is
\emph{witness} data, and a transaction's \texttt{txid} -- unlike its
\texttt{wtxid} -- is computed over the non-witness serialization only
(BIP-141). Two witness variants of the same underlying transaction can
therefore share an identical \texttt{txid} while carrying different
annexes, different Taproot signatures, and different \texttt{wtxid}s; only
one variant is ever mined, but a party who knows merely a \texttt{txid} has
no way to tell, from the \texttt{txid} alone, which witness/annex instance
that was. So a specific transaction must be identified by its
\texttt{wtxid}, or by some other authenticated source of the confirmed
witness -- a bare \texttt{txid} is not enough to pin down which annex is
even being talked about.

Closing the authentication gap itself decomposes into two separate
sub-properties, since a candidate $\mathrm{sha\_annex}$ and a signature
alone establish only the first:
\begin{enumerate}
\item \emph{candidate $\mathrm{sha\_annex} \to$ a validly signed witness
  variant} -- some witness instance exists for which this
  $\mathrm{sha\_annex}$, this signature, and these transaction/spent-output
  fields are mutually consistent;
\item \emph{that witness variant $\to$ actually confirmed on-chain} -- this
  specific witness instance, not some other validly-signed alternative the
  same signer could equally have produced, is the one Bitcoin's consensus
  actually accepted.
\end{enumerate}
Three candidate mechanisms, addressing these to different degrees:
\begin{itemize}
\item \textbf{Reconstruct $\mathit{SigMsg}$ and check the signature.}
  $\mathrm{sha\_annex}$ is one field embedded inside the larger
  $\mathit{SigMsg}$ that BIP-341's Taproot key-path signature is computed
  over (together with, among other fields, $\mathrm{sha\_prevouts}$,
  $\mathrm{sha\_amounts}$, $\mathrm{sha\_scriptpubkeys}$,
  $\mathrm{sha\_sequences}$, $\mathrm{sha\_outputs}$, and
  $\mathit{input\_index}$) -- it is \emph{not} independently serialized or
  recoverable from the final signature hash alone. A verifier who is given
  the full unsigned transaction (or the specific fields $\mathit{SigMsg}$
  needs), the spending input's Taproot public key, and its signature can
  recompute $\mathit{SigMsg}$ -- with the claimed $\mathrm{sha\_annex}$
  substituted in at the right position -- and check the BIP-340 Schnorr
  signature over $\mathrm{hash_{TapSighash}}(\mathtt{0x00} \Vert
  \mathit{SigMsg})$ directly. If it validates, the claimed
  $\mathrm{sha\_annex}$ is, up to breaking Schnorr unforgeability, part of
  a genuinely, validly signed witness for this transaction's inputs --
  \emph{sub-property 1 only}. It does \emph{not}, by itself, establish
  sub-property 2: the signature itself is witness data too, excluded from
  \texttt{txid} the same way the annex is, so a signer holding the private
  key can construct another equally valid signature over a different annex
  without changing the base transaction or its \texttt{txid} at all. This
  option proves authorization and internal consistency of a witness
  variant, not that this particular variant is the one Bitcoin confirmed.
\item \textbf{Trust a validating node or block explorer.} A lightweight
  verifier can simply ask a full node (or a service that has one) ``what is
  the annex hash for input $i$ of the transaction with this \texttt{wtxid}
  (or this block height and position)'', the same way it might already
  trust a block explorer for a confirmation status. Asking by
  \texttt{wtxid} (or an equivalent witness-committing identifier) rather
  than \texttt{txid} addresses sub-property 2 as well, at the cost of
  reintroducing a trusted third party the rest of this construction
  otherwise avoids.
\item \textbf{An authenticated copy of the confirmed witness, tied to the
  block's witness commitment.} BIP-141 commits every block's transactions'
  \texttt{wtxid}s into the coinbase transaction. A verifier who is given
  the witness commitment path from the claimed witness serialization up to
  a block header they otherwise trust (e.g.\ via an SPV-style proof of
  work) gets both sub-properties without trusting any single third party.
  But this route is \emph{not} $O(\log N)$: a \texttt{wtxid} is a double
  SHA-256 over the transaction's \emph{complete} serialization, witness
  included, and therefore over the annex bytes themselves. To check that a
  claimed $\mathrm{sha\_annex}$ belongs to the confirmed \texttt{wtxid}, the
  verifier must recompute the \texttt{wtxid} from the full witness, i.e.\
  hash all $\Theta(N)$ annex bytes. Nothing computed for one chain can be
  reused for the other: $\mathrm{sha\_annex}$ and the \texttt{wtxid} are two
  independent SHA-256 chains over the same bytes, started from different
  states, so FlatMerkle's checkpoints for the first say nothing about the
  second.
\end{itemize}
The three mechanisms therefore form a trilemma, and none of them gives a
lightweight verifier everything at once:
\begin{center}
\begin{tabular}{@{}lccc@{}}
\toprule
Mechanism & Confirmed variant & Trustless & Verifier cost \\
\midrule
Signature over $\mathit{SigMsg}$     & no  & yes & $O(1)$ \\
Trusted node / explorer              & yes & no  & $O(1)$ \\
Witness commitment to a header       & yes & yes & $O(N)$ \\
\bottomrule
\end{tabular}
\end{center}
The first option, on its own, reduces the scheme to an ordinary digital
signature over a commitment -- which needs no blockchain at all -- and so
gives up the timestamping and ordering guarantees listed in
Section~\ref{sec:annex-buys}; the second reintroduces exactly the trusted
party the construction is meant to avoid; the third is trustless and
confirmation-aware but costs as much as downloading the annex. Hence,
for a purely lightweight, trustless verifier, the $O(\log N)$ cost of a
FlatMerkle proof is the cost of checking membership \emph{given} an
authenticated $\mathrm{sha\_annex}$, not the end-to-end cost of learning
what Bitcoin confirmed. The $O(N)$ authentication step can be done once
per annex by any party and amortized across many proofs only if verifiers
are willing to trust that party's result, which is the second option
again. Settings with their own accountability structure can avoid the
trilemma: in the rollup of Section~\ref{sec:rollup-da}, sub-property~1 is
sufficient for the dispute itself, as explained there. None of these
mechanisms is implemented or benchmarked in this paper; a deployment needs
to pick one explicitly rather than assume FlatMerkle's own guarantees
cover it.

\subsection{Block alignment}
\label{sec:annex-alignment}

$\mathrm{sha\_annex} = \mathrm{SHA256}(\mathrm{compact\_size}(\mathrm{len}(\mathrm{annex})) \,\Vert\, \mathrm{annex})$
prepends a length preamble whose own size, $p \in \{1,3,5,9\}$ bytes,
depends on $\mathrm{len}(\mathrm{annex})$ (BIP-341's \texttt{CompactSize}
encoding). Since $p$ is essentially never a multiple of 64, treating this
as a simple byte-offset adjustment is not sufficient: unless corrected,
every 64-byte FlatMerkle item (Section~\ref{sec:construction}) straddles
two real SHA-256 compression-block boundaries, so no single item
corresponds to one
call of $\C$ at all, breaking the correspondence the whole construction
relies on. The fix is a small, fixed \emph{alignment header} of
$h = (64 - p) \bmod 64$ bytes -- since $p \in \{1,3,5,9\}$ is never a
multiple of 64, $h$ is always one of $\{63,61,59,55\}$, fixed once $p$ is
known (e.g.\ $h=59$ for the common case $p=5$) -- placed immediately after
the compact-size preamble and before the FlatMerkle structured file
proper, so that $p + h \equiv 0 \pmod{64}$: exactly one real SHA-256
compression block. Unlike an ordinary alignment padding scheme, this
header's content is \emph{not} free: BIP-341 recognizes a witness item as
an annex at all only if the annex's own first byte is \texttt{0x50}, so
the header's first byte is fixed to \texttt{0x50}. The remaining $h-1$
bytes carry no alignment requirement of their own -- padding \emph{length},
not content, is what restores block-alignment -- but this paper fixes them
canonically rather than leaving them to individual deployments: they
\texttt{SHALL} be all-zero bytes. Any other fixed convention (e.g.\ a magic
tag identifying the annex as a FlatMerkle structured file) would work
equally well for alignment, but leaving this genuinely a per-deployment
choice would mean two conforming FlatMerkle-in-annex implementations could
disagree on $\mathrm{effective}~\CK[0]$ for the byte-identical annex,
breaking interoperability; fixing it here removes that ambiguity.

With this header in place, define the \emph{effective} checkpoint
$\CK[0]$ (in place of the raw $\mathrm{IV}$, Section~\ref{sec:construction})
as the real SHA-256 chain state \emph{after} processing that one
preamble-plus-header block, starting from the ordinary $\mathrm{IV}$: every
FlatMerkle item from $D[0]$ onward then begins exactly on a genuine
64-byte compression-block boundary within the annex's own SHA-256
computation, and $\mathrm{position}(\cdot)$ (Lemma~\ref{lem:position})
converts directly to a byte offset by multiplying by 64 and adding
$p + h$, with no further adjustment needed. (A degenerate edge case is
worth flagging, not solving here: if $\mathrm{len}(\mathrm{annex})$ sits
exactly at one of the three \texttt{CompactSize} length thresholds -- 253,
65536, or $2^{32} = 4294967296$ bytes, where $p$ transitions from 1 to 3,
3 to 5, or 5 to 9 bytes respectively -- adding the header itself could
change $p$ and hence $h$; a real implementation should compute $p$ from
the \emph{final} annex length, header included, not iterate.)

Changing the starting chaining state alone is not enough to make this a
correct SHA-256 computation, however: SHA-256's finalization encodes the
\emph{total} number of bits processed from the true, original
$\mathrm{IV}$ -- not from whatever effective checkpoint a particular
application chooses to resume from. Standalone FlatMerkle's
$\mathrm{finalize}(\mathit{raw\_state}, N)$ (Algorithm~3.3) implicitly
assumes the entire hashed message is the $M = 2N-2$ structured-file items
themselves, $64M$ bytes total. In the annex, one additional real SHA-256
block -- the compact-size-preamble-plus-header block derived above -- was
hashed \emph{before} $\mathrm{effective}~\CK[0]$, so the true total message
length is $64(M+1)$ bytes, not $64M$: using the standalone
$\mathrm{finalize}(\cdot, N)$ unmodified would encode the wrong length and
simply fail to reproduce the real $\mathrm{sha\_annex}$, independent of
anything else being correct. The fix is to use an annex-specific
finalization that takes the extra prefix block count explicitly,
\begin{equation}
\mathrm{finalize\_annex}(\mathit{raw\_state}, N) := \mathrm{finalize}_{\,64(M+1)}(\mathit{raw\_state})
\end{equation}
where $\mathrm{finalize}_{\,L}(\cdot)$ denotes SHA-256's ordinary
length-suffix finalization parameterized directly by the true total byte
count $L$ (rather than inferring it solely from $N$), and
Algorithm~3.3's $\mathrm{Verify}$ must be instantiated with
$\mathrm{finalize\_annex}$, not the standalone $\mathrm{finalize}$, when
checking a proof against $\mathrm{sha\_annex}$.

\subsection{Caveats}

\begin{itemize}
\item \textbf{Relay and standardness policy.} As of this writing, annex
  data is consensus-valid but not relayed or mined by default nodes;
  alternative relay policies (e.g.\ Libre Relay, which relays annex-bearing
  transactions under a specific marker convention) exist but are not
  universal. Reliance on this application should account for the relay
  policy landscape at deployment time, which may change.
\item \textbf{Future consensus use of the annex.} BIP-341 currently defines
  the annex only as a reserved field for future extensions, and its own
  text is explicit and severe about this: ``[u]ntil the meaning of this
  field is defined by another softfork, users SHOULD NOT include annex in
  transactions, or it may lead to PERMANENT FUND LOSS'' \citep{bip341}.
  This is not a hypothetical or minor concern -- it is BIP-341's own
  strongest warning language, aimed at exactly this kind of use. It carries
  no consensus meaning today beyond being committed to via
  $\mathrm{sha\_annex}$, but a future soft fork could assign consensus-level
  semantics to some or all annex content (for instance, reserving a
  specific prefix byte or sub-format) that conflicts with using raw annex
  bytes for arbitrary application data as described here, in a way that
  could render funds in outputs spent through such a transaction
  unspendable. Anyone deploying this application should treat that risk as
  live, not theoretical, and track proposals that touch the annex field
  before relying on it.
\end{itemize}
